\chapter{Future Work} \label{ch:future_work}

While the methodology presented in this thesis successfully demonstrates the viability of intersecting g top-down system call graphs with bottom-up \textit{sysctl} knob graphs to identify performance-critical kernel configurations, it establishes a foundational framework that invites several promising directions for future work.

\section{Expanding the Syscall Analyser}

The current implementation of the Syscall Analyser was deliberately focused on the Linux networking subsystem to validate the core hypothesis.
A natural and necessary extension of this work is to broaden the analyser coverage to other critical kernel subsystems.
By mapping the execution paths of a wider variety of system calls (e.g., \texttt{mmap}, \texttt{madvise}, \texttt{fsync}, \texttt{clone}) the framework could be applied to entirely different classes of software, such as heavy relational databases or memory-intensive machine learning workloads.
Expanding the systemic vocabulary of the Syscall Analyser would allow it to isolate critical knobs in the Virtual File System (VFS), memory management, and process scheduling domains with the same precision currently achieved in the network stack.

\section{Mitigating the Curse of Dimensionality}

Although the graph intersection phase successfully filtered thousands of generic kernel parameters down to a targeted subset of 92 networking-specific knobs, this still represents a massive search space.
For automated tuning algorithms like Bayesian Optimization, exploring 92 dimensions simultaneously introduces a significant computational challenge known as the "curse of dimensionality".
To make automated tuning truly efficient, future iterations of this framework must further reduce this intersection set.
This dimensionality reduction can be directly achieved by addressing current methodological limitations through two primary directions: 
\begin{enumerate}
    \item \textbf{Enhancing static analysis precision}.
    \item \textbf{Incorporating dynamic runtime profiling}.
\end{enumerate}

\subsection{Enhancing Knob Analyser Precision}

A significant portion of the 92 identified knobs are false-positives generated by the inherent limitations of the current static analysis pipeline.
Currently, the Knob Analyser struggles to differentiate the precise execution paths of highly localized parameters, often generating nearly identical call graphs for all TCP knobs residing within core transmission and reception functions (e.g., \texttt{net.ipv4.tcp\_*}).

To prune these false positives and shrink the parameter space, the Knob Analyser must be enhanced with more granular, field-sensitive data flow tracking.
Rather than tracking the general propagation of a configuration structure, the analyser should be refined to track specific structure fields or array indices.
Furthermore, incorporating control-flow dependence analysis would allow the framework to determine not just if a knob value reaches a function, but whether the specific conditional branch utilizing that knob is actually logically reachable.

\subsection{Integration of Dynamic Analysis}

Even with perfect static analysis, static tools inherently suffer from over-approximation because they must account for all theoretically possible execution paths, including those hidden behind complex state machines or unresolvable function pointers.
To truly minimize the knob search space, future work should integrate dynamic analysis into the intersection pipeline.

By utilizing tracing tools like extended Berkeley Packet Filter (\textit{eBPF}) or \textit{ftrace,} the system could profile a real application under load to capture its exact, ground-truth execution path.
Intersecting this dynamic runtime graph with the static knob propagation graphs would create a highly precise hybrid analysis.
This hybrid approach would definitively filter out statically reachable but dynamically unused branches, cleanly eliminating irrelevant knobs and reducing the search space to the absolute minimum dimensions possible.

\section{Automated Reinforcement Learning Pipeline}

The ultimate application of this methodology is autonomous system tuning.
The benchmark environment utilized in this research, specifically the autotune framework, was also designed to generate structured datasets mapping kernel parameters to performance metrics.

A highly impactful future direction involves transitioning from offline, iterative Bayesian searches to an online Reinforcement Learning (RL) pipeline.
The mathematically filtered list of critical knobs generated by the methodology provides the exact, low-dimensional action space required to make RL feasible at the OS level.
By utilizing the datasets generated during this evaluation to pre-train an RL agent, future systems could deploy autonomous controllers capable of adjusting \textit{sysctl} parameters in real-time, instantly adapting to shifting workload bottlenecks without human intervention.

\newpage
